\section{Design Tic-Tac-Toe}
\label{sec:case-tic-tac-toe}

\problemstatement{Design a Tic-Tac-Toe game (generalisable to an $n\times
n$ board) supporting two players taking turns, move validation, and
win/draw detection. Keep the design extensible to other grid games.}

\begin{patternbox}[Clarify Before Designing]
Ask: Fixed $3\times3$ or general $n\times n$ with $k$-in-a-row to win? Two
human players or human-vs-AI? For this design assume: general board size,
$n$-in-a-row to win, two players, and an efficient $O(1)$ win check per
move rather than rescanning the board.
\end{patternbox}

\subsection*{Identify the entities}

\texttt{Game} (orchestrates turns and end conditions), \texttt{Board}
(grid state and win detection), \texttt{Player} (has a mark/symbol),
\texttt{Move}, and \texttt{Piece}/\texttt{Symbol}. The game is a small
lifecycle: in-progress, won, drawn -- a light State/enum.

\begin{ideabox}[Incremental Win Check in O(1) per Move]
Instead of scanning all lines after every move, keep running counts per
row, per column, and the two diagonals for each player. A move updates
just its row, column, and (if on a diagonal) the diagonal counters; a
count reaching $n$ means a win. This turns win detection from $O(n^2)$
into $O(1)$ per move.
\end{ideabox}

\subsection*{Class design (C++ skeleton)}

\begin{lstlisting}
#include <bits/stdc++.h>
using namespace std;

class Board {
    int n;
    vector<vector<char>> grid;
    vector<int> rowCnt, colCnt;                 // signed counts per player
    int diag = 0, anti = 0;
public:
    Board(int size) : n(size), grid(size, vector<char>(size, '.')),
                      rowCnt(size, 0), colCnt(size, 0) {}

    bool inBounds(int r, int c) const { return r >= 0 && r < n && c >= 0 && c < n; }
    bool isFree(int r, int c) const { return grid[r][c] == '.'; }

    // returns true if this move wins. player = +1 (X) or -1 (O)
    bool place(int r, int c, char mark, int player) {
        grid[r][c] = mark;
        rowCnt[r] += player; colCnt[c] += player;
        if (r == c) diag += player;
        if (r + c == n - 1) anti += player;
        int win = n * player;
        return rowCnt[r] == win || colCnt[c] == win || diag == win || anti == win;
    }
};

struct Player { string name; char mark; int id; };  // id = +1 or -1

class Game {
    Board board; vector<Player> players; int turn = 0; int movesLeft;
public:
    Game(int n, Player a, Player b)
        : board(n), players{a, b}, movesLeft(n * n) {}

    string move(int r, int c) {
        Player& p = players[turn];
        if (!board.inBounds(r, c) || !board.isFree(r, c)) return "Invalid move";
        bool won = board.place(r, c, p.mark, p.id);
        movesLeft--;
        if (won) return p.name + " wins";
        if (movesLeft == 0) return "Draw";
        turn ^= 1;                              // next player
        return "ok";
    }
};

int main() {
    Game g(3, {"Alice",'X',+1}, {"Bob",'O',-1});
    cout << g.move(0,0) << "\n";                // X
    cout << g.move(1,0) << "\n";                // O
    cout << g.move(0,1) << "\n";                // X
    cout << g.move(1,1) << "\n";                // O
    cout << g.move(0,2) << "\n";                // X wins (top row)
    return 0;
}
\end{lstlisting}

\begin{takeawaybox}
Tic-Tac-Toe rewards two things: a clean separation of \texttt{Game}
(rules/turns) from \texttt{Board} (state/win check), and the $O(1)$
incremental win-detection insight. Generalise to $n\times n$ up front.
Follow-ups: an AI player via \textbf{Strategy} (minimax), and other grid
games (Connect Four) reusing the board abstraction.
\end{takeawaybox}
